\documentclass[11pt]{article}

\usepackage[a4paper,margin=0.82in]{geometry}
\usepackage[T1]{fontenc}
\usepackage[utf8]{inputenc}
\usepackage{lmodern}
\usepackage{amsmath,amssymb,amsthm,mathtools}
\usepackage{enumitem,booktabs,array,tabularx,microtype,xcolor}
\usepackage[numbers,sort&compress]{natbib}
\usepackage[colorlinks=true,linkcolor=blue!65!black,
  citecolor=blue!65!black,urlcolor=blue!70!black]{hyperref}
\usepackage[nameinlink,noabbrev]{cleveref}
\setlength{\emergencystretch}{3em}

\hypersetup{
  pdftitle={A Source-Locked Census for the Production Phase Registry},
  pdfauthor={Liang Wang},
  pdfsubject={A missing-value census for the named physical phase and its determinant-two packet coordinates},
  pdfkeywords={production phase registry, named atom, fixed h0, packet coordinates, source lock}
}

\newtheorem{theorem}{Theorem}[section]
\newtheorem{proposition}[theorem]{Proposition}
\newtheorem{corollary}[theorem]{Corollary}
\theoremstyle{definition}
\newtheorem{definition}[theorem]{Definition}
\theoremstyle{remark}
\newtheorem{remark}[theorem]{Remark}

\newcommand{\T}{\mathbb T}
\newcommand{\NT}{\textnormal{\textsc{not-testable}}}
\newcommand{\Lzero}{\mathrm{L0}}
\newcommand{\Lone}{\mathrm{L1}}
\newcommand{\Ltwo}{\mathrm{L2}}
\newcolumntype{Y}{>{\raggedright\arraybackslash}X}
\newcolumntype{P}[1]{>{\raggedright\arraybackslash}p{#1}}

\title{\textbf{A Source-Locked Census for the Production Phase Registry:}\\
\textbf{Fixed-\(h_0=2\) Data, Missing Named-Atom Values,}\\
\textbf{and Determinant-Two Packet Coordinates}}
\author{Liang Wang\\
School of Artificial Intelligence and Automation,\\
Huazhong University of Science and Technology, Wuhan 430074, P.R. China\\
\texttt{wangliang.f@gmail.com}}
\date{July 2026}

\begin{document}
\maketitle

\begin{abstract}
The TPC-170 metric corridor controls all prefixes of every packet in
a prescribed determinant-two schedule for Lebesgue-almost every fixed
phase.  To consume that theorem at the physical endpoint one first
needs a value-bearing production phase registry: a source-located
named atom, its exact phase modulo one, and its map into the precise
packet coordinates on which the metric theorem acts.  We audit the
frozen TPC-157--172 phase-interface artifacts under canonical source
locks.  They prove that the relevant actual core has fixed
\(h_0=2\), and they specify the type of the missing H9 phase record.
They do not contain a value-bearing named atom, a phase-value locator,
a production packet schedule, or any physical-atom-to-packet row.

We give the minimal typed registry contract, including canonical
B\'ezout representative and covariant-multiplier fields, and return
\(\mathsf{NOT\_TESTABLE}\) with the first missing object equal to the
named atom together with its phase-value source locator.  The result
is a scoped \(\Lone\) source-interface obstruction.  It is not a
global nonexistence theorem, does not synthesize phase zero or any
other phase, and supplies no decay exponent.
\end{abstract}

\medskip
\noindent\fcolorbox{red!70!black}{red!3}{\parbox{0.94\linewidth}{
\textbf{Claim boundary.}
The fixed-\(h_0=2\) fact and the phase registry are separate data
objects.  An H9 registry has decay coordinate \(\mathsf{NONE}\).
Neither an obligation label, a source hash, a phase-average theorem,
nor a numerical fixture supplies a named physical phase.  No
fixed-atom estimate, program-positive \(\Ltwo\), strict \(1/400\),
prime-pair lower bound, or twin-prime theorem is proved.}}

\section{The exact question}

TPC-171 records the missing physical phase leaf as
\[
 \mathsf{H9.phase\_cell\_registry}.
\tag{1}\label{eq:h9}
\]
Its six-axis signature is
\[
\begin{array}{c|c}
\text{axis}&\text{required value}\\
\hline
\text{carrier}&\mathsf{PHYSICAL\_PHASE\_REGISTRY}\\
\text{phase}&\mathsf{NAMED\_FIXED\_ATOM}\\
\text{endpoint}&\mathsf{NOT\_APPLICABLE}\\
\text{scale}&\mathsf{DETERMINISTIC\_ALL\_SCALE}\\
\text{decay}&\mathsf{NONE}\\
\text{support}&\mathsf{ACTUAL\_ACTIVE\_SUPPORT}.
\end{array}
\tag{2}\label{eq:signature}
\]
The node has no arithmetic parent.  This direction is deliberate:
arithmetic theorems may consume a physical phase record, but cannot
manufacture one \citep{WangTPC161,WangTPC171}.

\begin{definition}[Value-bearing production phase record]
\label{def:record}
A production phase record contains at least:
\begin{enumerate}[label=(\roman*),leftmargin=2.2em]
\item a stable named physical atom identifier;
\item its exact value \(\alpha_\star\in\T\);
\item a theorem-usable source path and locator for that value;
\item a separate source lock for \(h_0=2\);
\item the exact production packet schedule to which the phase is
  mapped; and
\item source-located rows realizing the map into the packet
  coordinates of TPC-170.
\end{enumerate}
\end{definition}

An object that contains only the words ``named phase'', only a target
node such as \eqref{eq:h9}, or only a schema allowing these fields is
not a record in the sense of \cref{def:record}.

\section{\texorpdfstring{Fixed \(h_0=2\)}{Fixed h0=2} is present and is not a phase}

TPC-157 and TPC-159 both export the determinant-two actual-core field
\[
 h_0=2
\tag{3}\label{eq:h0}
\]
in their machine-readable theorem records
\citep{WangTPC157,WangTPC159}.  Endpoint Ledger V4 in TPC-171 likewise
marks the fixed physical \(h_0\) gate as proved on the relevant actual
core \citep{WangTPC171}.  These three locators are retained
independently in the census.

\begin{proposition}[Data-axis separation]
\label{prop:separation}
The source-backed fact \eqref{eq:h0} does not instantiate any of the
first three fields of \cref{def:record}.
\end{proposition}

\begin{proof}
The integer \(h_0\) is the prescribed physical shift, while
\(\alpha_\star\) is a character coordinate in
\(\T=\mathbb R/\mathbb Z\).  Their types and source fields differ.
The cited records contain the integer \(2\) in a
\texttt{fixed\_h0} or fixed-physical-\(h_0\) field.  None gives a
named-atom identifier, a phase value modulo one, or a locator for such
a value.  Reusing the integer \(2\), reducing it modulo one, or
silently declaring \(\alpha_\star=0\) would be a type-changing
invention.
\end{proof}

This proposition does not weaken the fixed-\(h_0\) result.  It only
prevents that genuine data fact from filling a different missing
field.

\section{Packet coordinates required by TPC-170}

For a determinant-two packet, TPC-170 uses
\[
 su-ad=2,\qquad q=as,\qquad
 t_d(z)=ad+qz,
\tag{4}\label{eq:fiber}
\]
and the coefficient sequence
\[
 c_{d,u}(z)=\mu(d+sz)\mu(u+az).
\tag{5}\label{eq:coeff}
\]
At ambient scale \(X_n\), a prescribed finite list
\(\mathcal P_n\) contains a canonical fiber, a terminal scale
\(T_{n,p}\), and a multiplier \(\rho_{n,p}\).  Its ordered
coordinates \(z_{n,p,j}\) enter
\[
 S_{n,p,k}(\alpha)
 =
 \sum_{j\le k}
 c_{d,u}(z_{n,p,j})\rho_{n,p}(z_{n,p,j})
 e(-\alpha z_{n,p,j}).
\tag{6}\label{eq:sum}
\]
Thus a production crosswalk must give actual values for
\[
\begin{split}
(&X_n,p,a,s,d,u,q,T_{n,p},\rho_{n,p},z_{n,p,j}),\\
&\text{the physical occurrence locator, and the common }
\alpha_\star.
\end{split}
\tag{7}\label{eq:coordinates}
\]
The formula \eqref{eq:sum} defines a consumer interface; it is not a
production row.

There is also a representative condition.  For fixed \((a,s)\), all
solutions of \eqref{eq:fiber} are
\[
 (d,u)\longmapsto(d+sr,u+ar),\qquad r\in\mathbb Z.
\tag{8}\label{eq:representative}
\]
TPC-170 proves that this translates the coefficient packet.  The
multiplier must be translated covariantly; then the twist changes by
only a unit scalar.  A production registry must therefore select one
canonical representative or record the translation index and the
covariant multiplier map \citep{WangTPC170}.  The theorem
\eqref{eq:representative} does not select the actual production
representative by itself.

\section{Frozen source census}

The executable census locks the phase-interface outputs of
TPC-157--161 and TPC-167--172.  Its relevant findings are summarized
in \cref{tab:census}.

\begin{table}[ht]
\centering
\small
\begin{tabularx}{\textwidth}{P{0.29\textwidth}P{0.25\textwidth}Y}
\toprule
Object & Frozen value & Exact meaning\\
\midrule
fixed physical \(h_0\) & \(2\), source-backed &
One literal data gate is present.\\
TPC-158 production phase cell & \(\NT\) &
The major/minor interface has no production phase value.\\
TPC-161 phase artifact & \(\mathsf{MISSING}\) &
The target description is present; the required artifact is not.\\
TPC-171 H9 phase node & \(\NT\), evidence null &
The six-axis obligation is typed but uninstantiated.\\
TPC-170 phase mode & Lebesgue-a.e.\ fixed phase &
A metric quantifier, not a named atom.\\
named atom/value/locator rows & \(0/0/0\) &
No value-bearing production phase record occurs in the seven
explicitly mapped fields.\\
production packet-coordinate rows & \(0\) &
The mapped theorem/status fields supply no row of
\eqref{eq:coordinates}.\\
\bottomrule
\end{tabularx}
\caption{The fixed-\(h_0\) fact survives the census; the production
phase values do not appear.}
\label{tab:census}
\end{table}

TPC-167 proves a continuous phase-\(L^2\) identity, TPC-168 a
separated-registry density theorem, TPC-169 a maximal prefix theorem,
and TPC-170 an almost-every-fixed-phase corridor
\citep{WangTPC167,WangTPC168,WangTPC169,WangTPC170}.  Their claim
boundaries all reject a named production phase.  They are positive
arithmetic inputs in their declared norms, not value-bearing entries
for \eqref{eq:h9}.

\begin{theorem}[Scoped registry census]
\label{thm:census}
In the frozen source-locked TPC-157--172 phase-interface corpus:
\begin{enumerate}[label=(\roman*),leftmargin=2.2em]
\item \(h_0=2\) is source-backed on the relevant actual core;
\item the target type \eqref{eq:signature} is source-backed;
\item the representative rule \eqref{eq:representative} is proved;
\item none of the seven explicitly mapped fields is a value-bearing
  named physical atom or phase-value locator; and
\item those mapped fields contain no production row realizing
  \eqref{eq:coordinates}.
\end{enumerate}
Consequently the production phase registry has status \(\NT\).
\end{theorem}

\begin{proof}
The first three statements are exact field checks against the
source-locked TPC-157, TPC-159, TPC-170 and TPC-171 outputs.  The last
two statements are an explicit mapped-field census: three
fixed-\(h_0\) records and four phase-obligation records.  It is not a
generic scanner, and newly added source fields require an explicit
mapping update.  The machine audit rejects any snapshot that fills a
null phase value, locator, schedule, or packet row without a new
source.  By \cref{def:record}, the mapped evidence does not construct
a production registry.
\end{proof}

\begin{remark}[Scope]
\label{rem:scope}
\Cref{thm:census} is not a theorem that no physical phase description
exists mathematically or in every possible archive.  It says that the
explicitly mapped fields in the declared frozen corpus do not provide
the values needed to instantiate \cref{def:record}.  A larger corpus
or a new field must be named, mapped, locked, and audited before it can
change the result.
\end{remark}

\section{Decision and typed progress}

The exact machine result is
\[
\boxed{
\begin{gathered}
\mathsf{Verdict}=\mathsf{NOT\_TESTABLE},\\
\mathsf{ProductionPhaseRegistryConstructed}=\mathsf{false},\\
\mathsf{FirstMissing}=
\mathsf{named\_physical\_atom\_id\_and\_}\\[-2pt]
\mathsf{phase\_value\_source\_locator}.
\end{gathered}}
\tag{9}\label{eq:decision}
\]
Independent missing fields include the production schedule locator,
the packet-coordinate rows, and the actual representative/covariant
multiplier map.

The progress levels are:
\begin{itemize}[leftmargin=2em]
\item \(\Lzero\): deterministic hashes, schemas, field census, and
  mutation diagnostics;
\item \(\Lone\): the scoped source-interface conclusion
  \cref{thm:census}; and
\item \(\Ltwo\): none.
\end{itemize}
The fixed-\(h_0=2\) field is a proved literal data fact, but it creates
no cancellation exponent.  The H9 phase leaf retains decay axis
\(\mathsf{NONE}\).

\section{Reproducible artifacts and next gate}

The generator writes
\begin{verbatim}
experiments/tpc180_phase_registry_census.json
experiments/tpc180_phase_registry_census_audit.json
\end{verbatim}
and checks both strict schemas.  Source hashes have integrity
semantics only.  Mutation regressions reject a target obligation used
as a phase value, an unsourced locator, a synthetic packet row, H9
decay promotion, a scoped census promoted to global nonexistence, and
the fixed-\(h_0\) fact promoted to \(\Ltwo\).

The next arithmetic object remains
\[
 \mathsf{H2.metric\_fixed\_atom\_crosswalk}.
\tag{10}\label{eq:next}
\]
Even a future positive phase registry would supply data only.  TPC-181
must separately ask whether a theorem places its named atom in the
TPC-170 good set.  No such theorem is inferred here.

\section{Conclusion}

The phase question is now split at the correct seam.  Fixed
\(h_0=2\) and the determinant-two Fourier interface are present; the
value-bearing named phase and its production packet crosswalk are not.
The only source-compatible TPC-180 output is therefore the contract
and scoped missing-value census \eqref{eq:decision}.  This is useful
\(\Lone\) localization, not fixed-atom arithmetic progress.

\nocite{*}
\bibliographystyle{plainnat}
\bibliography{references}

\end{document}
